\section{Formalization}\label{sec:formal}

\begin{figure}[t!]
  \vspace{-.3cm}
{\footnotesize
  \begin{alignat*}{2}
    \text{\textbf{Variables }}& \quad &\quad& x, y, z, \vnu, ...
    \\\text{\textbf{Base Types}}& \quad & b  ::= \quad &  \Code{unit} ~|~ \Code{bool} ~|~ \Code{nat} ~|~ \Code{int} ~|~ ...
    \\\text{\textbf{Pure Operations}}& \quad & \primop ::= \quad & {+} ~|~ {-} ~|~ {==} ~|~ {<} ~|~ {\leq} ~|~ ...
    \\ \text{\textbf{Constants }}& \quad &c ::= \quad & () ~|~ \mathbb{B} ~|~ \mathbb{Z} ~|~ ...
    \\ \text{\textbf{Values }}& \quad  & v ::= \quad & c ~|~ x 
    \\ \text{\textbf{Qualifiers}}& \quad & \phi ::= \quad & v ~|~ \primop\;\overline{v} ~|~ \bot ~|~ \top  ~|~ \neg \phi ~|~ \phi \land \phi ~|~ \phi \lor \phi ~|~ \phi \impl \phi ~|~ \forall x{:}b.\, \phi
    \\\text{\textbf{Effectful Operations}}& \quad & \eff{op} ::= \quad & \eff{readReq} ~|~ \eff{readRsp} ~|~  ...
    \\\text{\textbf{ event Kinds}}& \quad & k ::= \quad & \eff{gen} ~|~ \eff{obs}
    \\ \text{\textbf{Expressions}}& \quad & e ::=\quad & v ~|~ \zlet{x{:}b}{\primop\ \overline{v}}{e} ~|~ \zobs{\overline{x{:}b}}{op}{e} ~|~ \zlet{\overline{x{:}b}}{\sassume{\phi}}{e}
    \\ & & \quad & |~ \zgen{op}{\overline{v}}{e} ~|~ \zassert{\phi}{e} ~|~ \zloop{e}{e} ~|~ e \oplus e
  \end{alignat*}}
\vspace*{-.25in}
\caption{$\DSL$ syntax }
\label{fig:term-syntax}
\vspace*{-.1in}
\end{figure}

We formalize our approach using a core language, $\DSL$, for
expressing controller programs. This language is an asynchronous event-driven calculus in A-normal form (ANF)\cite{FSDF93} that abstracts away the
implementation details of the library that the controller interacts
with, focusing only on the structure of the controller program
itself. The syntax of $\DSL$ is shown in \autoref{fig:term-syntax}; 
it includes both pure and effectful operations ($\primop$ and
$\eff{op}$), non-deterministic choice and loops ($\oplus$ and $\mykw{loop}$), and assertions.
Effectful operations are categorized as either
\emph{generable} ($\eff{gen}$) or \emph{observable} ($\eff{obs}$).


\begin{figure}[t!]
\vspace*{-0.1in}
{\footnotesize
{\small
\begin{flalign*}
\text{\textbf{Events}} \ m ::= \  \zevent{op}{\overline{c}} &
\parbox{7cm}{\centering$\text{\textbf{capability}}\  \beta \in \mathcal{P}(m)$}  &
\text{\textbf{Traces}} \ \alpha ::= \ \emptr ~|~ m \;{::}\; \alpha ~|~ \alpha \listconcat \alpha &
\end{flalign*}
\vspace*{-.2in}
\begin{flalign*}
  &\text{\textbf{Handler Semantics }}\ 
  \fbox{$\alpha \vDash \zevent{op}{\overline{c}} \Downarrow \beta$} &
  \text{\textbf{Operational Semantics }}\ 
  \fbox{$\phi \Downarrow c \quad \steptr{\alpha}{(\buf, e)}{\alpha}{(\buf, e)}$}
\end{flalign*}}
    \begin{prooftree}
    \mhypo{50mm}
    {$\buf = \{\zevent{op}{\overline{c}}\} \cup \buf_1 $ \quad
      $\alpha \vDash \zevent{op}{\overline{c}} \Downarrow \buf_2$ \quad
      $e' = e[\overline{x \mapsto c}]$
    }
      \infer1[\textsc{\small StObs}]{
        \steptr{\alpha}{(\buf,
          \zobs{\overline{x}}{op}{e})}{[\zevent{op}{\overline{c}}]}{(\buf_1
          \cup \buf_2, e')}
      }
    \end{prooftree}
      \quad
    \begin{prooftree}
      \hypo{\alpha \vDash \zevent{op}{\overline{c}} \Downarrow \beta' }
      \infer1[\textsc{\small StGen}]{
        \steptr{\alpha}{(\buf, \zgen{op}{\overline{c}}{e})}{[\zevent{op}{\overline{c}}]}{(\buf \cup \beta', e)}
      }
    \end{prooftree}
\\ \
\mprooftr{58mm}{
$\phi\msubst{x}{c}\Downarrow \top$
}{StAssume}{
$\steptr{\alpha}{(\buf, \zlet{\overline{x}}{\sassume{\phi}}{e})}{[]}{(\buf, e\msubst{x}{c})}$
}
\quad
\mprooftr{44mm}{$\phi \Downarrow \top$}{StAssert}{
$\steptr{\alpha}{(\buf, \zassert{\phi}{e})}{[]}{(\buf, e)}$
}
% \\[0.4em]
% \mprooftr{45mm}{}{StLoop1}{
% $\steptr{\alpha}{(\buf, \zloop{e_1}{e_2})}{[]}{(\buf, e_2)}$
% }
% \quad
% \mprooftr{70mm}{}{StLoop2}{
% $\steptr{\alpha}{(\buf, \zloop{e_1}{e_2})}{[]}{(\buf, e_1;\zloop{e_1}{e_2})}$
% }
}
\vspace*{-.1in}
  \caption{Selected Operational Semantics\ZZ{I still call it capability. The operational semantics allow a send request doesn't has a response -- however, our typing rule guarentee it is really linearly consumed. }}
  \label{fig:selected-semantics}
\vspace*{-.2in}
\end{figure}

\paragraph{Operational Semantics} Events in $\DSL$ are operations
applied to concrete values $(\zevent{op}{\overline{c}})$. Evaluating a
$\DSL$ program depends on an context \emph{trace}, i.e., a sequence of
events, and an \emph{observation capability}, i.e., a multiset of
observable events. Each evaluation step produces an output trace and an updated
capability. Traces are equipped with the standard list operations (i.e.,
cons ${::}$ and concatenation $\listconcat{}$).  The operational
semantics of $\DSL$ are defined by the small-step reduction relation:
$\steptr{\alpha}{(\beta, e)}{~\alpha'}{(\beta', e')}$.  This judgment
is read as: ``under the context trace $\alpha$ and observation capability
$\beta$, $e$ steps to $e'$, emitting the trace $\alpha'$ and remaining
the capability $\beta'$.''  Intuitively, the context trace $\alpha$
represents the sequence of events visible to a handler, thereby
determining its response; the capability $\beta$ contains events that
are able to be observed in the future. The semantics uses an auxiliary judgement,
$\alpha \vDash \zevent{op}{\overline{c}} \Downarrow \beta$, that
specifies any new events that need to be added to the observation capability
after handling $\eff{op}$.

\autoref{fig:selected-semantics} provides the key rules of $\DSL$'s
semantics.\footnote{The remaining rules are completely standard and
  provided in the supplemental material.} The rule for observable
events (\textsc{StObs}) reflects the ``consume-and-produce'' semantics of
observation capability. This rule non-deterministically removes
a event in capability that matches the effectful operation $\eff{op}$,
evaluates it under the current context, and substitutes the event
payload $\overline{c}$ for the variables $\overline{x}$ in $e$, the
body of the $\mykw{let}$ expression. Any new observable events generated as a
consequence of handling $\eff{op}\ \overline{c}$ are added to the
resulting  observation capability.  The reduction rule for generable events
(\textsc{StGen}) is similar, but since these events can be directly
performed by the controller, the rule does not require a corresponding
 event in the capability. The \textsc{StAssume} rule substitutes the
variables $\overline{x}$ with values $\overline{c}$ that satisfy the
qualifier $\phi$ in the body of an assume expression. The
\textsc{StAssert} rule, in contrast, requires the qualifier of an
assert expression to hold in order for it to make progress.

\ZZ{TODO: Rewrite this example after rewrite Sec 2}
\begin{example}[Operational Semantics]\label{ex:semantics}
  The first three events in the trace \ref{eqn:traceEx} are produced by
  the controller program \ref{eqn:autoEx} as follows:
  {\footnotesize
    \begin{align*}
      {[]}\vDash{(\emptyset,P_C)}
      &\xhookrightarrow{[]} (\emptyset,\emph{lines $2$ - $7$ of $P_C$, with $x\mapsto 3, y \mapsto 4$}) \tag{\textsc{StAssume}}\\[-0.4em]
      &\xhookrightarrow{[ \zevent{writeReq}{3}]} (\{ \zevent{writeRsp}{3}\},
        \emph{lines $3$ - $7$ of $P_C$}) \tag{\textsc{StGen}}\\[-0.4em]
      &\xhookrightarrow{[\zevent{writeReq}{4}]} (\{\zevent{writeRsp}{3},
        \zevent{writeRsp}{4}\}, \emph{lines $4$ - $7$ of $P_C$}) \tag{\textsc{StGen}}\\[-0.4em]
      &\xhookrightarrow{[\zevent{writeRsp}{4}]} (\{\zevent{writeRsp}{3}\},
        \emph{lines $5$ - $7$ of $P_C$}) \tag{\textsc{StObs, StAssert}} \\[-0.4em]
      &\xhookrightarrow{[\eff{readReq}]} (\{\zevent{writeRsp}{3}, \zevent{readRsp}{4,\true}\},
        \emph{lines $6$ - $7$ of $P_C$}) \tag{\textsc{StGen}}
    \end{align*}
  }\noindent The first step performs the substitution
  ($x\mapsto 3, y \mapsto 4$), which satisfies the assumed formula
  $x\not=y$ (line $1$). In the next two steps, $P_{C}$ generates two
  $\eff{writeReq}$  events and adds two $\eff{writeRsp}$  events to
  the  observation capability. One of these  events is consumed by the fourth
  step, causing the assertion on line 4 of \autoref{fig:2pc-term} to
  succeed. The fifth step handles $\eff{readReq}$, and the  event
  $\eff{readRsp(4, true)}$ is added to the capability.
\end{example}

\subsection{Types}

\begin{figure}[t!]
{\small
\begin{alignat*}{2}
    \text{\textbf{Pure Refinement Types}}& \quad & t  ::= \quad & \rawnuot{b}{\phi} ~|~ x{:}t \sarr t
    \\\text{\textbf{Symbolic Event}}& \quad &se  ::= \quad & \msgB{op}{\overline{x}}{\phi}
    \\\text{\textbf{Symbolic Automata}}& \quad &H,A,F  ::= \quad & \text{\textcolor{gray}{in Symbolic Regex }} \emptyset ~|~ \epsilon ~|~ se ~|~ A \lorA A ~|~ A \seqA A ~|~ A^* ~|~ \negA A ~|~ A \landA A ~|~ \anyA
    \\& \quad &  \quad & \text{\textcolor{gray}{in Symbolic LTL$_f$ }}  se ~|~ A \lorA A ~|~ \nextA A ~|~ A \untilA A ~|~ \negA A ~|~ A \landA A 
    \\ \text{\textbf{Prophecy Automata Types}}& \quad & \tau  ::= \quad &
    \rg{H}{A}{F} ~|~ x{:}b \garr \tau ~|~ x{:}t\sarr \tau ~|~ \tau \interty \tau
    \\\text{\textbf{Type Contexts}}& \quad &\Gamma ::= \quad  &\emptyset ~|~ x{:}t, \Gamma
  \end{alignat*}
  \vspace*{-.1in}
  \begin{flalign*}
    \text{\textbf{Type Alias}}& \quad \msg{op}{\phi} \doteq \msgB{op}{\overline{x}}{\phi} \text{ when the feild name $\overline{x}$ are clear} \quad  \ha{H}{A} \doteq \rg{H}{A}{\allA}
  \end{flalign*}
}
\vspace*{-.15in}
\caption{\langname{} types.}
\label{fig:type-syntax}
\vspace*{-.15in}
\end{figure}


The syntax of types in $\DSL{}$ is shown in \autoref{fig:type-syntax}.
Types include \emph{pure} refinement types, which describe pure
computations, and Prophecy Automata Types (\PAT{}s), which describe
effectful computations.  Pure refinement types are similar to those
found in other refinement type systems~\cite{JV21}, and allow base
types (e.g., {\Code{int}}) to be further constrained by a logical
formula or qualifier. Verification conditions generated by our
type-checker can be encoded as effectively propositional (EPR)
sentences~\cite{Ramsey1987}, which can be efficiently handled by an
off-the-shelf theorem prover such as Z3~\cite{de2008z3}.

\paragraph{Symbolic Finite Automata} Following other recent
trace-based type systems\cite{ZYDJ24}, $\DSL{}$ uses \emph{Symbolic
  Finite Automata} (SFAs)~\cite{Vea13,SFA-minimization,
  SFA-Transducers} to qualify traces, similar to how standard
refinement types use formulae to qualify the types of pure terms.
We represent SFAs using a symbolic version of regular expressions; symbolic LTL$_f$ is also supported and can be translated to regex as needed.
A \emph{symbolic event} $\msgB{op}{\overline{x}}{\phi}$ is an atomic predicate that
describes an effectful operation {$\eff{op}$} whose inputs
$\overline{x}$ must satisfy the qualifier $\phi$.
% \footnote{When the
%   fields of an event are clear from context, we omit its parameters
%   $\overline{x}$, e.g., $\msg{writeReq}{v > 0}$ means
%   $\msgB{writeReq}{v}{v > 0}$.} 
  The standard regex operators (i.e., empty set of traces $\emptyset$, empty trace $\epsilon$, union $A \lorA A$, sequence $A \seqA A$, and kleene star $A^*$) and various set operators (i.e., complement $\neg$, intersection $\land$, and an union of all symbolic events $\anyA$) are defined normally. When the fields of a symbolic event are clear from context, we omit its parameters $\overline{x}$, e.g., $\msg{writeReq}{v > 0}$ means $\msgB{writeReq}{v}{v > 0}$.
  
% The standard temporal operators
% (e.g., test $\evparenth{\phi}$, next $\nextA A$, until $\untilA$) and
% various set operators (i.e., complement $\neg$, intersection $\land$,
% and union $\lorA$) are defined normally. These operators are
% sufficient to capture other modalities, e.g., eventually ($\finalA$),
% globally ($\globalA$), and importantly, the singleton (last) modality
% $\lastA$, which describes a singleton trace, i.e., one which prohibits
% any subsequent effects~\cite{LTLf}. SFAs can capture several common
% patterns: the set of all possible traces \(\allTL\), the singleton set
% containing the empty trace \(\epsilon\), and the empty set of traces and
% \(\empTL\); these are analogous to the regular expressions \( .^* \),
% \(\epsilon\), and \(\emptyset\), resp. 



\paragraph{Prophecy Automata Types} A \PAT{} $\rg{H}{A}{F}$ is
comprised of three SFAs: a \emph{history} SFA $H$ that captures the
context traces (i.e., a sequence of visible, already handled, symbolic
events) in which a term can be executed, a \emph{current} SFA $A$ that
describes newly handled  events that arise from executing a term, and
a \emph{prophecy} SFA $F$ that characterizes events that have to be performed in the future. When a \PAT has no contraint on future computation, we omit its prophecy automata as $\ha{H}{A}$. Function types use \PAT{}s in their result types
to describe the effects they perform, when combined with intersection
types ($\interty$), this allows users to express complex control
flows. Function types also use \emph{ghost variables}
($x{:}b \garr \tau$) to capture data dependencies among
symbolic events; for example, the full signature of the $\eff{getReq}$
handler from \autoref{sec:intro} uses the ghost variables $\Code{key}$
and $\Code{val}$.

\ZZ{TODO: rewrite this example after rewrite Sec 2}
\begin{example}[Strong Consistency]
  Strong consistency requires that all $\eff{getRsp}$  events report
  the last value that was $\eff{put}$ to the database. This property
  is captured by the following \PAT{}:
%
  {
    \footnotesize
    \begin{align*}
      \Code{val}\hasnty{tVal} \garr & \Code{key}\hasoty{tKey}{\top} \sarr  \effLR{\finalA \msg{putReq}{k = \Code{key} \land v = \Code{val}} \land \nextA \neg \finalA \msg{putReq}{k = \Code{key}} } %
      \\& \quad \effLR{\msg{getReq}{k = \Code{key}}}%
      \effLR{(\neg \msg{putReq}{k = \Code{key}}) \untilA \msg{getRsp}{k = \Code{key} \land v = \Code{val}}}
    \end{align*}}\noindent
  The prophecy automata in this \PAT{} requires that no
  updates ($\eff{putReq}$) to $\Code{key}$ in the database
  happen before a user receives a response to a $\eff{getReq}$  event
  for the key $\Code{key}$.
\end{example}

\subsection{Typing rules}
Our typing judgment features three contexts: a type context $\Gamma$,
a handler context $\Delta$, and a capability context $\Theta$. The
type context, $\Gamma$ maps from variables to pure refinement types
(i.e., $t$). As in other trace-based refinement type systems, contexts
are not allowed to contain \PAT{}s--- doing so breaks several
structural properties (e.g., weakening) that are used to prove type
safety. The handler context, $\Delta$, provide the relation between operations, specification of its handler as a \PAT{}, and the operations
its handler adds to the capability. The capability context, $\Theta$,
records the multiset of operations could be observed. This context
is used to ensure that every observation corresponds to a event that
was triggered by a previous event. The type judgement guarantee all capability in $\Theta$ are consumed by the term linearly and it has no contraint on future computation, i.e., the term's type always has the form $\ha{H}{A}$.

\begin{example} The handler context $\Delta$ for our running examples
  augments the four specifications from \autoref{fig:spec} as follows:
{\small
\begin{align*} 
  \Delta \doteq \{ (\eff{readReq},  \{\eff{readRsp}\}, \ldots),%
      (\eff{readRsp}, \emptyset, \ldots),%
      (\eff{writeReq}, \{\eff{writeRsp}\},\ldots), %
      (\eff{writeRsp}, \emptyset, \ldots) \}
\end{align*}
}
\end{example}\vspace{-.2cm}

\begin{figure}[!t]
{\footnotesize
{\small
\begin{flalign*}
 &\text{\textbf{Auxiliary Typing}} & \fbox{$\Gamma \wellfoundedvdash \tau \quad \Gamma \vdash A \subseteq A \quad \Gamma \vdash \tau <: \tau$} &
  &\text{\textbf{Typing}} &  &\fbox{$\Gamma \vdash v : t \quad \Gamma;~ \Delta;~ \Theta \vdash e: \tau$}
\end{flalign*}
}
\\ \
% \mprooftr{38mm}
% {$\Gamma \wellfoundedvdash H$ \quad
% $\Gamma \wellfoundedvdash A$ \quad
% $\Gamma \wellfoundedvdash F$ }
% {WfHAF}
% {$\Gamma \wellfoundedvdash \rg{H}{A}{F}$}
% \
\mprooftr{42mm}
{$\Gamma \vdash H_2 \subseteq H_1$ \quad 
$\Gamma \vdash F_2 \subseteq F_1$ \quad
$\Gamma \vdash A_1 \subseteq A_2$  }
{SubHAF}
{$\Gamma \vdash \rg{H_1}{A_1}{F_1} <: \rg{H_2}{A_2}{F_2}$}
\
\mprooftr{20mm}
{$\Gamma; \Delta; \Theta \vdash e: \tau$ \quad
$\Gamma; \Delta; \Theta \vdash \tau <: \tau'$}
{TSub}
{$\Gamma; \Delta; \Theta \vdash e: \tau'$}
\\[.4em]
\mprooftr{21mm}
{
$\Gamma; \Delta; \Theta \vdash e_1 : \tau$ \quad $\Gamma; \Delta; \Theta \vdash e_2 : \tau$
 }
{TChoice}
{$\Gamma; \Delta; \Theta \vdash e_1 \oplus e_2 : \tau$}
\quad
\mprooftr{87mm}
{
$(\eff{op}, \Theta', \overline{x_i{:}t_i}\sarr
\rg{H}{\msg{op}{\phi}}{A\seqA \allA}) \in \Delta$ \quad
$\forall i. \Gamma \vdash v_i : t_i$ \quad
 $\Gamma; \Delta; \Theta \cup \Theta' \vdash e : \ha{H \seqA \msg{op}{\phi\msubst{x_i}{v_i}}}{A}$
 }
{TGen}
{$\Gamma; \Delta; \Theta \vdash \zgen{op}{\overline{v_i}}{e} : \ha{H}{\msg{op}{\phi\msubst{x_i}{v_i}} \seqA A}$}
\\[.4em]
\mprooftr{35mm}
{}
{TRet}
{$\Gamma; \Delta; \emptyset \vdash () : \ha{\allA}{\epsilon}$}
\quad
\mprooftr{81mm}
{
  $(\eff{op}, \Theta', \overline{x_i{:}t_i}\sarr \rg{H}{\msgB{op}{\overline{y}}{\phi}}{A\seqA \allA}) \in \Delta$ \quad
 $\Gamma, \overline{x{:}t}; \Delta; \Theta \cup \Theta' \vdash e : \ha{H \seqA \lastA\msgB{op}{\overline{y}}{\phi \land \overline{y = x}}}{A}$
 }
{TObs}
{$\Gamma; \Delta; \{\eff{op}\} \cup \Theta \vdash \zobs{\overline{x}}{op}{e} : \ha{H}{\lastA\msgB{op}{\overline{y}}{\phi} \seqA A}$}
\\[.4em]
\mprooftr{61mm}{
 $\Gamma; \Delta; \emptyset \vdash e_1 : \ha{H \seqA A_1^*}{A_1}$ \quad
 $\Gamma; \Delta; \Theta \vdash e_2 : \ha{H \seqA A_1^*}{A_2}$
}{\textcolor{red}{TLoop}}{
 $\Gamma; \Delta; \Theta \vdash \zloop{e_1}{e_2} : \ha{H}{A_1^*\seqA A_2}$}
}
% \vspace*{-.1in}
\caption{Selected typing rules.}
\label{fig:selected-typing-rules}
\vspace*{-.15in}
\end{figure}

\paragraph{Auxiliary typing relations}
Our system depends on three auxiliary relations: a well-formedness
relation $\Gamma \wellfoundedvdash \tau$ which ensures, e.g., that all
qualifiers appearing in a type $\tau$ are closed under the current
typing context $\Gamma$; an inclusion relation on SFAs
$\Gamma \vdash A \subseteq A$; and a mostly-standard semantic
subtyping relation.  \autoref{fig:selected-typing-rules} provides one
of the key rules for these relations.\ZZ{Well-formedness rule is not interesting to introduce since we don't require non-emptyness in type level} 
% A well-formed \PAT{}
% ($\textsc{WfHAF}$) is required to accept at least one trace
% ($\neg\allTL$ is an SFA that rejects all traces). 
Subtyping for two
\PAT{}s ($\textsc{SubHAF}$) is established by checking inclusion
between their constituent automata under the current type context
$\Gamma$. Inclusion on the history and prophecy automata is
contravariant, while current automata are covariant. 
The prophecy automata specify constraints on possible future event interleavings, restricting only subset of future computations that can satisfy the type, which is same as how the history automata constrain the context trace. In contrast, the current automata overapproximate all traces that the term itself may produce.
Intuitively, since both the history and prophecy automata restrict the contexts in which a term that produces the current automata may appear, it is safe to further constrain both contexts.


% \paragraph{Effect Operator Framing}  The use of prophecy
% automata in specifying handler behavior is very powerful.  However,
% since it is a characterization of future events that must happen as a
% consequence of executing a handler  event, it is important that the
% specification does not overly constrain the effects that may be
% produced by other handlers that execute independently.  In other
% words, a well-formed handler specification (defined in the next
% subsection) should have its prophecy automata be in the form
% $\bigwedge_{i = 1 ... n} \finalA \msg{op_i}{\phi_i}$.  In order to
% allow this type to adapt to independently occuring events in the
% future, we introduce a framing rule \textsc{TEOpFrame} for operators
% that allows the arbitrary insertion of other automata ($F_1
% ... F_{n+1}$) between $\msg{op_i}{\phi_i}$.

\paragraph{Typing Rules} A subset of our typing rules is shown in
\autoref{fig:selected-typing-rules}.\footnote{The complete set of
  typing rules is included in the \techreport{}.} All of our terms
assume any types they use are well-formed, so we elide the
corresponding well-formedness judgments from their premises. 
The rules for performing events, \textsc{TGen} and \textsc{TObs}, both extract
the type of the corresponding handler from $\Delta$,
$\rg{H}{\msg{op}{\phi}}{A}$, and require that it aligns
with the \PAT{} of the expression that the operation is being
performed in:%
{\footnotesize
  \begin{align*}
    \text{\textcolor{gray}{(term's \PAT)}}
    \underbrace{ H
    }_\texttt{history}\;\seqA\;\underbrace{
    \msg{op}{\phi\msubst{x_i}{v_i}} \seqA A }_\texttt{current}
    \;\seqA\; \underbrace{ \allA }_\texttt{prophecy} \equiv \underbrace{ H
    }_\texttt{history}\;\seqA\;\underbrace{
    \msg{op}{\phi\msubst{x_i}{v_i}} }_\texttt{current} \;\seqA\;
    \underbrace{ A \seqA \allA}_\texttt{prophecy}
    \text{\textcolor{gray}{(operator's \PAT)}}
  \end{align*}}\noindent %
To type the rest of the expression, both rules move the symbolic event
$\msg{op}{\phi\msubst{x_i}{v_i}}$ from the head of the current
automata to the tail of the history automata and add any new
capabilities to $\Theta$. In order to make an observation on
$\eff{op}$, \textsc{TObs} additionally requires that the capability
context has a corresponding capability ($\{\eff{op}\} \cup \Theta$). The
standard subsumption rule \textsc{TSub} allows us to change the shape
of a type that a term is being typed against. Controllers always end in a
unit value $()$; thus, the \textsc{TRet} rule requires the current
automata of this term ($\epsilon$) to only accept the empty trace (i.e.,
$[]$). The nondeterministic choice operator is typed using the
\textsc{TChoice} rule, when combined with the subsumption rule, this
allows controllers to explore different event orderings. 
The nondeterministic loop aligns with the Kleene star in SFA, and thus can be straightforwardly typed by the \textsc{TLoop} rule.
To preserve linearity, the loop body is prohibited from consuming events from the capability context $\Theta$, since the number of iterations is unbounded.

% of observing effect operator $\eff{op}$ from the typing
% context with the help of effect operator typing
% $\textsc{TEOp}$. Besides requiring the arguments ($\overline{v_i}$) is
% well-typed against parameter types ($\overline{t_i}$) as application
% rules in standard refinement typing, it additionally

\ZZ{TODO: rewrite this example after rewrite Sec 2}
\begin{example}[Controller Typing] We provide an informal typing
  derivation of $P_C$ against a \PAT{} that encodes a violation of an RYR
  policy, $\rg{\epsilon}{\negSCA}{\epsilon}$, under the type context
  $\Gamma \doteq \Code{x}\hasoty{int}{\top}, \Code{y}\hasoty{int}{v
    \not= x}$. The first step of our derivation uses \textsc{TSub} to
  refine our target type to a \PAT{} that better aligns with the
   events sent by $P_C$: %
  {\footnotesize
    \begin{align*}
      \negSCA' \doteq
      &\underbrace{
        \lastA\msg{writeReq}{v = \Code{x}}
        }_\texttt{$A_1$} \seqA \underbrace{ \lastA\msg{writeReq}{v =
        \Code{y} \land v \not= \Code{x}} \seqA
        \lastA\msg{writeRsp}{v = \Code{y}} \seqA
        \lastA\msg{readReq}{\top}
        }_\texttt{$A_2$} \seqA \\[-.3em]
      &\underbrace{ \lastA\msg{writeRsp}{v =
        \Code{x}}
        }_\texttt{$A_3$} \seqA \underbrace{ \lastA\msg{readRsp}{v =
        \Code{y} \land st = \true \land v \not= \Code{x}}
        }_\texttt{$A_4$}
    \end{align*}}\noindent
  The first $\mykw{gen}$ expression on line $2$ of $P_C$
  is then typed using \textsc{TGen}. After retrieving the specification of
  $\eff{writeReq}$ from $\Delta$ and uses \textsc{TSub} to adjust it
  into a shape consistent with $\negSCA'$:
  {\footnotesize
    \begin{align*}
      &\Delta(\eff{writeReq}) = \mykw{gen}\;\langle \Code{x}\hasnty{int}
        \sarr \rg{\allTL}{\lastA\msg{writeReq}{v =
        \Code{x}}}{\finalA\msg{writeRsp}{v = \Code{x}}}, \{\eff{writeRsp}\} \rangle \\
      & \Gamma \vdash \Code{x}\hasnty{int}
        \sarr \rg{\allTL}{\lastA\msg{writeReq}{v =
        \Code{x}}}{\finalA\msg{writeRsp}{v = \Code{x}}} \\
        & \underline{\;\quad<: \Code{x}\hasnty{int} \sarr \effLR{\epsilon}\
        \effLR{\lastA\msg{writeReq}{v = \Code{x}}}\ \effLR{A_2 \seqA
        \lastA\msg{writeRsp}{v = \Code{x}} \seqA A_4}} \\
      &\Gamma;\Delta; \emptyset \vdash P_{C} \text{ (lines $2$ - $7$)}: \rg{\epsilon}{A_1\seqA A_2\seqA A_3\seqA A_4}{\epsilon}
        \tag{\textsc{TGen}}
    \end{align*}
    \begin{align*}
      &(\eff{writeReq}, \Code{x}\hasnty{int} \sarr \rg{\epsilon}{\msg{writeReq}{v =\Code{x}}}{A_2\seqA\msg{writeRsp}{v = \Code{x}}\seqA A_4}, \{\eff{writeRsp}\}) \in \Delta \\
      & \underline{\Gamma;\Delta; \emptyset \vdash P_{C} \text{ (lines $3$ - $7$)}: \ha{A_1}{A_2\seqA A_3\seqA A_4}} \\
      &\Gamma;\Delta; \emptyset \vdash P_{C} \text{ (lines $2$ - $7$)}: \ha{\epsilon}{A_1\seqA A_2\seqA A_3\seqA A_4}
        \tag{\textsc{TGen}}
    \end{align*}
  }\noindent %
Since the type of $\eff{writeReq}$ aligns with the
  target type $\rg{\epsilon}{A_1\seqA A_2\seqA A_3\seqA A_4}{\epsilon}$,
  we continue typing the rest of $P_C$ (lines $3$ - $7$) against the
  \PAT{} $\rg{A_1}{A_2\seqA A_3\seqA A_4}{\epsilon}$.
\end{example}

\subsection{Type Soundness}

\begin{figure}[t!]
  {\small\begin{align*}
    \denot{\rg{H}{A}{F}} \doteq &\{e ~|~ \emptyset \basicvdash e : \Code{unit} \land \forall \alpha_h \in\denot{H}.\forall \alpha_f \in\denot{F}. \forall \alpha\; \beta_1\;\beta_2\;e_h\;e_f.
    \\& \ \underbrace{\msteptr{[]}{(\emptyset, e_h)}{\alpha_h}{(\beta_1, ())}}_\texttt{realizable histroy trace} \land \underbrace{\msteptr{\alpha_h}{(\beta_1, e)}{\alpha}{(\beta_2, ())}}_\texttt{trace produced by $e$} \land \underbrace{\msteptr{\alpha_h \listconcat \alpha }{(\beta_2, e_f)}{\alpha_f}{(\emptyset, ())}}_\texttt{realizable future trace} \impl \alpha \in\denot{A} \}
  \end{align*}}
  \caption{Selected type denotation.}
  \label{fig:selected-denotation}
  \vspace*{-.15in}
  \end{figure}

\paragraph{Type denotations} Similar to other refinement type systems~\cite{JV21},
types in $\DSL$ are denoted as their inhabitants (i.e., $\denot{t}$ and $\denot{\tau}$).  The capability context is denoted as  observation capabilitys, while the type context $\Gamma$ is denoted as \emph{substitution} $\sigma$ (e.g., $[x\mapsto 3, y\mapsto 4]$ in Example~\ref{ex:semantics}) that provides the assignments for binding variables in $\Gamma$. Moreover, the denotation (accepting language)
of SFAs is the set of traces they can accept. Then,  automata
inclusion under a type context is defined as $\Gamma \vdash A \subseteq
A' \doteq \forall \sigma \in \denot{\Gamma}. \denot{\sigma(A)}
\subseteq \denot{\sigma(B)}$. 
\autoref{fig:selected-denotation} defines the denotation of $\rg{H}{A}{F}$: a term $e$ inhabits $\rg{H}{A}{F}$ iff, for all realizable (i.e., producible by the execution of some program) history traces $\alpha_h$ and future traces $\alpha_f$ accepted by $H$ and $F$, any current trace $\alpha$ produced by $e$ is accepted by $A$.
\footnote{The details of these definition can be found in our \techreport{}.}

\paragraph{Handler specification} 
We expect the trace produced by our test generator to be consistent with the handler specification $\Delta$. That is, from the point of view of each event in the trace, the history, current, and future traces are consistent with the corresponding $\PAT$ in $\Delta$. On the other hand, $\Delta$ should be consistent with the
auxiliary semantics of handlers introduced in \autoref{fig:selected-semantics}.

\begin{definition}[Trace consistent with handler context]
  \label{lemma:trace-hctx}
  A trace $\alpha$ is consistent with a handler context $\Delta$ (written $\traceDelta{\Delta}{\alpha}$) iff for every $\zevent{op}{\overline{c_i}}$ in $\alpha$ (i.e., $\alpha = \alpha_h \listconcat [\zevent{op}{\overline{c_i}}] \listconcat \alpha_f$), there exists $\overline{x{:}t} \sarr \rg{H}{\msg{op}{\phi}}{F}$ and capability $\Theta$ in $\Delta$ such that the following hold:
{\small\begin{align*}
    \alpha_h \in \denot{H\msubst{x_i}{c_i}} \land \alpha_f \in \denot{F\msubst{x_i}{c_i}} \land \Theta \subseteq \{ m.\eff{op} ~|~ m \in \alpha_f\} \land \phi\msubst{x_i}{c_i} \land \forall i. c_i \in \denotation{t_i}
\end{align*}}
\end{definition}

% \paragraph{Well-formed Handler specification} A
% handler specification $\Delta$ should be consistent with the
% auxiliary semantics of handlers introduced in
% \autoref{fig:selected-semantics}, also, $\Delta$ should also guarantee the new sending  event assumed by capability context can be eventually received.

\begin{definition}[Well-formed handler context]
  \label{lemma:built-in-typing}
  The handler specification $\Delta$ is well-formed iff for every $(\eff{op}, \Theta, \tau) \in \Delta$, the type $\tau$ is of the form $\overline{y{:}b} \garr \overline{x{:}t} \sarr \rg{H}{\msg{op}{\phi}}{F}$ and the following holds:
{\small\begin{align*}
    &\forall \overline{y{:}b}.\forall \overline{c \in \denotation{t}}.\forall \alpha_h \in \denot{H}. \forall \overline{m_i}. \alpha_h \vDash \zevent{op}{\overline{c}} \Downarrow \{\overline{m_i}\} \impl 
    \\&\quad (\forall \alpha_f \in \denot{F}. \alpha_f = \alpha_1 \listconcat [m_1] \listconcat ... [m_n] \listconcat \alpha_{n+1}) \land \phi\msubst{x}{c} \land \Theta = \{\overline{m_i.\eff{op}}\}
\end{align*}}
\end{definition}



\begin{theorem}[Fundamental Theorem]\label{theorem:fundamental}A well-typed term, i.e., \(\Gamma; \Delta; \Theta \vdash e : \ha{H}{A}\), generates traces consistent with the \PAT: $\forall \sigma \in \denotation{\Gamma}. \sigma(e) \in \denotation{\sigma(\ha{H}{A})}$.\footnote{The proofs of all theorems in this paper are provided in the \techreport{}.}
\end{theorem}
  
\begin{corollary}[Type Soundness]\label{theorem:sound} Given a well-formed handler specification \(\Delta\), with ghost variables \(\overline{x{:}b}\) and a violation property \(A\), a controller \(e\) that satisfies \(\overline{x{:\rawnuot{b}{\top}}};\Delta;\emptyset \vdash e : \ha{\epsilon}{A}\) will consistent with both violation property and handler context, i.e., 
   {\small\begin{align*} \forall
      \overline{c{:}b}. \forall \alpha. \msteptr{[]}{(\emptyset,
        e[\overline{x\mapsto c}])}{\alpha}{(\emptyset, ())} \implies \alpha \in
      \denot{A[\overline{x\mapsto c}]} \land \traceDelta{\Delta}{\alpha}
  \end{align*}}
\end{corollary}